\clearpage
\begin{summarybox}
\begin{center}\begin{tikzpicture}[font=\small,>=stealth,line join=round]
\draw[gray!40,rounded corners=3pt] (0,0) rectangle (14.2,5.8);\node at (7.1,5.2) {考前一眼：真数正，过(1,0)，大底升、小底降};
\node at (3.1999999999999997,4.55) {$a>1$：递增};\begin{scope}[shift={(1.2209302325581395,2.9)},x=1.069767441860465cm,y=0.35714285714285715cm]\draw[->,gray] (-.3,0)--(4,0) node[right]{$x$};\draw[->,gray] (0,-2.8)--(0,2.8) node[above]{$y$};\draw[blue!75!black,thick,domain=.22:3.7,samples=150] plot(\x,{ln(\x)/ln(2)});\fill[orange!85!black] (1,0) circle(2pt);\end{scope}\node at (3.1999999999999997,1.4) {负\quad$|\ 0\ |$\quad 正};\node at (10.0,4.55) {$0<a<1$：递减};\begin{scope}[shift={(8.02093023255814,2.9)},x=1.069767441860465cm,y=0.35714285714285715cm]\draw[->,gray] (-.3,0)--(4,0) node[right]{$x$};\draw[->,gray] (0,-2.8)--(0,2.8) node[above]{$y$};\draw[orange!85!black,thick,domain=.22:3.7,samples=150] plot(\x,{ln(\x)/ln(0.5)});\fill[orange!85!black] (1,0) circle(2pt);\end{scope}\node at (10.0,1.4) {正\quad$|\ 0\ |$\quad 负};
\node[text width=12.6cm,align=center] at (7.1,.65) {从左到右读符号；分界真数为1。图分别取底数2与1/2，其余合法底数按同类判断。};
\end{tikzpicture}\end{center}



	\textbf{\textcolor{CSummary}{一句话核心}}

	对数是指数运算的“倒过来问”：$a^y=x\iff y=\log_a x$。真数必须为正，函数值可正可负；恒过$(1,0)$；$a>1$时递增，$0<a<1$时递减。

	\medskip
	\textbf{\textcolor{CSummary}{使用条件}}
	\begin{itemize}[leftmargin=*, itemsep=0.5em, topsep=0.4em]
		\item \textbf{条件1（真数为正）：} 真数必须满足$x>0$。
		\item \textbf{条件2（底数合规）：} 底数$a>0$且$a\neq 1$。
	\end{itemize}

	\medskip
	\textbf{\textcolor{CSummary}{关键提醒}}
	\begin{itemize}[leftmargin=*, itemsep=0.5em, topsep=0.4em]
		\item \textbf{易错点（定义域优先）：} 任何对数问题首先检查真数大于$0$。
		\item \textbf{检查项（底数分类）：} 底数含参数时，必须分$a>1$和$0<a<1$讨论。
\end{itemize}

\medskip
\textbf{\textcolor{CSummary}{快记（每条都能解释）}}
\begin{itemize}[leftmargin=*,itemsep=0.35em,topsep=0.3em]
 \item \textbf{真数正：}因为$a^y$始终为正；不是说$\log_a x$始终为正。
 \item \textbf{过$(1,0)$：}因为$a^0=1$，所以$\log_a1=0$。
 \item \textbf{大底升、小底降：}指数$2^y$随$y$升高而变大，$(\tfrac12)^y$则变小；倒过来问“要几次方”时，方向对应不变。
 \item \textbf{两种对称：}底数互为倒数的两条对数图像关于$x$轴；指数与对应的对数图像关于$y=x$。
\end{itemize}

\medskip
\textbf{\textcolor{CSummary}{快用（三步速答）}}

\begin{center}
\begin{tikzpicture}[>=stealth]
 \node[draw=gray!55,rounded corners,inner sep=5pt] (D) at (0,0) {\small 真数$>0$};
 \node[draw=gray!55,rounded corners,inner sep=5pt] (A) at (3.0,0) {\small 看底数$a$};
 \node[draw=gray!55,rounded corners,inner sep=5pt] (R) at (6.15,0) {\small 比真数、验边界};
 \draw[->,thick] (D)--(A); \draw[->,thick] (A)--(R);
\end{tikzpicture}
\end{center}

\textbf{答题时机：}同底对数比较、含参对数不等式、图像识别、求固定经过点，优先用这三步，不用先算小数。不同底数难比较时，先试$0$或$1$作中间值。
\end{summarybox}
